\scanpage{83}
Pour montrer que $u$ est induit par un
autom. (unique) de $M$, il est tentant
d'écrire $M$, en termes de \deleted{l'action}
\deleted{de} $(\widetilde S_M,\widetilde F_M,T_\bullet:\Lambda^*\,-\,\unsure{\circlearrowright|\mathfrak S}_{\widetilde S_M})$, comme
\[
\begin{aligned}
M\xleftarrow{\sim}{}&\text{quotient du \inserted{$\mathbb Z$-}module libre }\mathbb Z^{(M^*)}\\
&\text{par le sous-groupe engendré}\\
&\text{par les }\ [x+y]-[x]-[y]\\
&\text{(pour }x,y,\deleted{\unsure{x+y}}\in M^*\ \text{tels que }x\wedge y\in(\det M)^*)\\
&\text{et par les }\ [-x]+[x]\;.
\end{aligned}\tag*{(23)}
\]
% slanted
\margin{\uline{NB} La structure ``triangulation'' $+$ ``antipod.'' peut s'interpréter (en termes de l'axiome de l'unicité du triangle complétant une arête) par une loi de groupe comm. pas partout définie $(x,y)\mapsto x\unsure{\mathbin{\perp}|\mathbin{+}}y$, $x\mapsto -x$, dont $M$ pourrait alors s'interpréter comme la ``\unsure{quasi-addition|quasi-loi|groupe abélien|groupe}'' \unsure{enveloppante|enveloppant}\ldots}

Désignant par $R_M$ ce sous-groupe,
une suite d'homomorphismes can.
\diagram{\(\text{(24)}\quad\begin{tikzcd}[ampersand replacement=\&,column sep=1.3em]0\arrow[r]\&R_M\arrow[r,"i"]\&\mathbb Z^{M^*}\arrow[r,"p"]\&M\arrow[r]\&0\end{tikzcd}\)}
et il faudrait montrer qu'elle est
exacte (le seul pb \unsure{ici|il} est $\operatorname{Ker}p=\operatorname{Im}i$,
alors qu'il est clair que $pi=0$). (Une
question analogue peut se poser pour un
module libre de rang quelconque \inserted{$\geq 2$}.)

Ainsi, $u$ satisfaisant (21) induit une
\deleted{application} additif $\widehat u$ de $M$ qui l'induit,
et la condition (22) implique que ce
dernier commute à $\Lambda^*$, donc à $\Lambda$ si $\Lambda^*$
engendre $\Lambda$ additivement\ldots{} Il est clair a
% archive sheet number: 82
